\documentclass[preview]{standalone}

\usepackage{amsmath}
\usepackage{amssymb}
\usepackage{tikz}
\usepackage{stellar}
\usepackage{definitions}
\usepackage{bettelini}

\begin{document}

\id{linearalgebra-dimension}
\genpage

\section{Dimension}

\begin{snippetproposition}{basis-have-same-cardinality}{Basis cardinality}
    Let \(V\) be a finite-dimensional \vectorspace over a \field
    \(\mathbb{F}\). Any two \basis[bases] of \(V\) have the same cardinality.
\end{snippetproposition}

\begin{snippetproof}{basis-have-same-cardinality-proof}{basis-have-same-cardinality}{Basis cardinality}
    Let
    \[
        B=\{v_1,\ldots,v_n\},
        \qquad
        C=\{w_1,\ldots,w_m\}
    \]
    be two \basis[bases] of \(V\). Since \(B\) is \linearlyindependent and \(C\)
    generates \(V\), the comparison theorem between independent collections
    and generating collections gives \(n\leq m\). Since \(C\) is
    \linearlyindependent and \(B\) generates \(V\), the same theorem gives
    \(m\leq n\). Hence \(m=n\).
\end{snippetproof}

\begin{snippetdefinition}{vector-space-dimension-definition}{Vector space dimension}
    Let \(V\) be a finite-dimensional \vectorspace. The \emph{dimension} of
    \(V\), denoted
    \[
        \lineardim(V),
    \]
    is the cardinality of any \basis of \(V\).
\end{snippetdefinition}

\begin{snippetexample}{fundamental-dimensions-example}{Fundamental dimensions}
    Let \(\mathbb{F}\) be a \field. Then
    \[
        \lineardim(\mathbb{F}^n)=n,
    \]
    because the canonical \basis has \(n\) elements. Also,
    \[
        \lineardim(\{0\})=0,
    \]
    because the zero \vectorspace has the empty \basis. For bounded-degree
    \polynomial[polynomials],
    \[
        \lineardim\mathcal{P}_m(\mathbb{F})=m+1,
    \]
    because \(\{1,x,\ldots,x^m\}\) is a \basis. Finally,
    \[
        \lineardim\matrices_{m\times n}(\mathbb{F})=mn,
    \]
    because the elementary \matrix[matrices] \(E_{ij}\) form a \basis.
\end{snippetexample}

\begin{snippetnote}{dimension-depends-on-scalar-field}{Dimension depends on the scalar field}
    If the scalar \field is not clear from context, it should be shown in the
    notation. For example,
    \[
        \lineardim_{\complexnumbers}\complexnumbers=1,
        \qquad
        \lineardim_{\realnumbers}\complexnumbers=2.
    \]
    As a complex \vectorspace, \(\complexnumbers\) has \basis \(\{1\}\). As a
    real \vectorspace, it has \basis \(\{1,\mathrm{i}\}\).
\end{snippetnote}

\section{Subspaces}

\begin{snippetproposition}{linear-subspace-has-lower-dimension}{Dimension of a subspace}
    Let \(V\) be a finite-dimensional \vectorspace over a \field
    \(\mathbb{F}\), and let \(U\subseteq V\) be a linear subspace. Then
    \[
        \lineardim(U)\leq\lineardim(V).
    \]
\end{snippetproposition}

\begin{snippetproof}{linear-subspace-has-lower-dimension-proof}{linear-subspace-has-lower-dimension}{Dimension of a subspace}
    A \basis of \(U\) is \linearlyindependent as a collection of \vector[vectors] of
    \(V\). A \basis of \(V\) generates \(V\). By the comparison theorem
    between independent collections and generating collections, a \basis of
    \(U\) has at most as many elements as a \basis of \(V\).
\end{snippetproof}

\begin{snippetproposition}{same-dimension-subspace-equality}{Subspaces with full dimension}
    Let \(V\) be a finite-dimensional \vectorspace over a \field
    \(\mathbb{F}\), and let \(U\subseteq V\) be a linear subspace. Then
    \[
        \lineardim(U)=\lineardim(V)
        \qquad
        \Longleftrightarrow
        \qquad
        U=V.
    \]
\end{snippetproposition}

\begin{snippetproof}{same-dimension-subspace-equality-proof}{same-dimension-subspace-equality}{Subspaces with full dimension}
    \begin{iffproof}
        \item Suppose that \(\lineardim(U)=\lineardim(V)=n\). Let
        \(\{u_1,\ldots,u_n\}\) be a \basis of \(U\). Since \(U\subseteq V\),
        this collection is \linearlyindependent in \(V\). By the \basis
        extension theorem, it can be extended to a \basis of \(V\). However,
        it already has \(n\) elements, so no \vector can be added. Therefore
        it is a \basis of \(V\), and
        \[
            V=\linearspan(u_1,\ldots,u_n)=U.
        \]

        \item If \(U=V\), then clearly \(\lineardim(U)=\lineardim(V)\).
    \end{iffproof}
\end{snippetproof}

\section{Dimensional criteria}

\begin{snippetproposition}{linearly-independent-set-of-dim-is-basis}{Independent collection of full size}
    Let \(V\) be a finite-dimensional \vectorspace over a \field
    \(\mathbb{F}\), with \(\lineardim(V)=n\). If
    \[
        \{v_1,\ldots,v_n\}
    \]
    is \linearlyindependent, then it is a \basis of \(V\).
\end{snippetproposition}

\begin{snippetproof}{linearly-independent-set-of-dim-is-basis-proof}{linearly-independent-set-of-dim-is-basis}{Independent collection of full size}
    A \linearlyindependent collection can be extended to a \basis of \(V\).
    Every \basis of \(V\) has exactly \(n\) elements, and the given collection
    already has \(n\) elements. Hence no additional \vector can be added, so
    the given collection is already a \basis.
\end{snippetproof}

\begin{snippetproposition}{generating-set-of-dim-is-basis}{Generating collection of full size}
    Let \(V\) be a finite-dimensional \vectorspace over a \field
    \(\mathbb{F}\), with \(\lineardim(V)=n\). If
    \[
        \{v_1,\ldots,v_n\}
    \]
    generates \(V\), then it is a \basis of \(V\).
\end{snippetproposition}

\begin{snippetproof}{generating-set-of-dim-is-basis-proof}{generating-set-of-dim-is-basis}{Generating collection of full size}
    A finite generating collection can be reduced to a \basis by removing
    redundant \vector[vectors]. Every \basis of \(V\) has exactly \(n\) elements, and
    the given generating collection already has \(n\) elements. Therefore no
    \vector can be removed, and the collection is already a \basis.
\end{snippetproof}

\section{Grassmann formula}

\begin{snippet}{grassmann-formula-sets-analogy}
    For finite \set[sets],
    \[
        |A\union B|
        =
        |A|+|B|-|A\intersection B|.
    \]
    Grassmann's formula is the \vectorspace analogue: \lineardimtext[dimensions] replace
    cardinalities, and sums of subspaces replace unions.
\end{snippet}

\begin{snippettheorem}{grassman-formula-theorem}{Grassmann formula}
    Let \(V\) be a finite-dimensional \vectorspace over a \field
    \(\mathbb{F}\), and let \(U,W\) be linear subspaces of \(V\). Then
    \[
        \lineardim(U+W)
        =
        \lineardim(U)+\lineardim(W)-\lineardim(U\intersection W).
    \]
\end{snippettheorem}

\begin{snippetproof}{grassman-formula-theorem-proof}{grassman-formula-theorem}{Grassmann formula}
    Since \(V\) is finite-dimensional, the subspaces \(U\), \(W\),
    \(U\intersection W\), and \(U+W\) are finite-dimensional. Let
    \[
        \{v_1,\ldots,v_h\}
    \]
    be a \basis of \(U\intersection W\). Extend it to a \basis of \(U\):
    \[
        \{v_1,\ldots,v_h,u_1,\ldots,u_j\},
    \]
    and extend it to a \basis of \(W\):
    \[
        \{v_1,\ldots,v_h,w_1,\ldots,w_k\}.
    \]
    We show that
    \[
        \mathcal{B}
        =
        \{v_1,\ldots,v_h,u_1,\ldots,u_j,w_1,\ldots,w_k\}
    \]
    is a \basis of \(U+W\).

    First, \(\mathcal{B}\) generates \(U+W\). If \(x\in U+W\), then
    \(x=u+w\) for some \(u\in U\) and \(w\in W\). Expressing \(u\) in the
    extended \basis of \(U\) and \(w\) in the extended \basis of \(W\), the sum
    \(u+w\) is a \linearcombination of \vector[vectors] in \(\mathcal{B}\).

    Now suppose that
    \[
        a_1v_1+\cdots+a_hv_h
        +b_1u_1+\cdots+b_ju_j
        +c_1w_1+\cdots+c_kw_k
        =
        0.
    \]
    Then
    \[
        c_1w_1+\cdots+c_kw_k
        =
        -(a_1v_1+\cdots+a_hv_h+b_1u_1+\cdots+b_ju_j).
    \]
    The left-hand side belongs to \(W\), and the right-hand side belongs to
    \(U\), so this \vector belongs to \(U\intersection W\). Since
    \(\{v_1,\ldots,v_h\}\) is a \basis of \(U\intersection W\), there are
    \vsscalar[scalars] \(d_1,\ldots,d_h\) such that
    \[
        c_1w_1+\cdots+c_kw_k=d_1v_1+\cdots+d_hv_h.
    \]
    Hence
    \[
        d_1v_1+\cdots+d_hv_h-c_1w_1-\cdots-c_kw_k=0.
    \]
    Since \(\{v_1,\ldots,v_h,w_1,\ldots,w_k\}\) is a \basis of \(W\), all
    these coefficients are zero. Thus \(c_1=\cdots=c_k=0\). The original
    relation becomes
    \[
        a_1v_1+\cdots+a_hv_h+b_1u_1+\cdots+b_ju_j=0.
    \]
    Since \(\{v_1,\ldots,v_h,u_1,\ldots,u_j\}\) is a \basis of \(U\), all
    remaining coefficients are zero. Therefore \(\mathcal{B}\) is
    \linearlyindependent.

    Hence \(\mathcal{B}\) is a \basis of \(U+W\), so
    \[
        \lineardim(U+W)=h+j+k.
    \]
    Since
    \[
        \lineardim(U)=h+j,
        \qquad
        \lineardim(W)=h+k,
        \qquad
        \lineardim(U\intersection W)=h,
    \]
    the formula follows.
\end{snippetproof}

\section{Examples}

\begin{snippetexample}{grassmann-two-planes-r3-example}{Two planes in \(\mathbb{R}^3\)}
    Let
    \[
        \Pi_1=\{(x,y,z)\in\realnumbers^3\suchthat x+y+z=0\},
    \]
    and
    \[
        \Pi_2=\{(x,y,z)\in\realnumbers^3\suchthat x-y+2z=0\}.
    \]
    Both are linear subspaces of \(\realnumbers^3\). We have
    \[
        \Pi_1=\linearspan((1,0,-1),(0,1,-1)),
        \qquad
        \lineardim(\Pi_1)=2,
    \]
    and
    \[
        \Pi_2=\linearspan((1,1,0),(0,2,1)),
        \qquad
        \lineardim(\Pi_2)=2.
    \]
    The intersection is determined by
    \[
        \begin{cases}
            x+y+z=0,\\
            x-y+2z=0.
        \end{cases}
    \]
    Solving gives
    \[
        (x,y,z)=y(-3,1,2),
    \]
    so
    \[
        \Pi_1\intersection\Pi_2=\linearspan((-3,1,2)),
        \qquad
        \lineardim(\Pi_1\intersection\Pi_2)=1.
    \]
    Therefore, by Grassmann's formula,
    \[
        \lineardim(\Pi_1+\Pi_2)=2+2-1=3.
    \]
    Since \(\Pi_1+\Pi_2\subseteq\realnumbers^3\) and both have \lineardimtext
    \(3\), it follows that
    \[
        \Pi_1+\Pi_2=\realnumbers^3.
    \]
\end{snippetexample}

\begin{snippetexample}{grassmann-two-five-dimensional-subspaces-r9-example}{Two five-dimensional subspaces of \(\mathbb{R}^9\)}
    Let \(U,W\subseteq\realnumbers^9\) be linear subspaces such that
    \[
        \lineardim(U)=5,
        \qquad
        \lineardim(W)=5.
    \]
    Since \(U+W\subseteq\realnumbers^9\),
    \[
        \lineardim(U+W)\leq 9.
    \]
    Grassmann's formula gives
    \[
        \lineardim(U\intersection W)
        =
        \lineardim(U)+\lineardim(W)-\lineardim(U+W)
        \geq
        5+5-9
        =
        1.
    \]
    Thus \(U\intersection W\neq\{0\}\). In particular, two
    five-dimensional subspaces of \(\realnumbers^9\) always intersect in at
    least a vector line.
\end{snippetexample}

\end{document}
